\documentclass[a4paper,11pt]{article}
\usepackage[utf8]{inputenc}
\usepackage[T1]{fontenc}
\usepackage[french]{babel}
\usepackage{lmodern}
\usepackage{amsmath,amssymb,amsthm}
\usepackage{mathrsfs}
\usepackage{enumitem}
\usepackage{multicol}

\setlist[enumerate]{label=\textbf{\textcolor{Couleur8}{\arabic*.}}, left=1em}

% Si vous voulez aussi modifier les listes enumerate imbriquées :
\setlist[enumerate,2]{label=\textcolor{Couleur8}{\alph*)}}
\setlist[enumerate,3]{label=\textcolor{Couleur8}{\roman*)}}
\setlist[enumerate,4]{label=\textcolor{Couleur8}{\Alph*)}}
\usepackage{graphicx}
\usepackage{xcolor}
\usepackage{tcolorbox}
\usepackage{geometry}
\usepackage{fancyhdr}
\usepackage{titlesec}
\usepackage{cancel}
\usepackage{ifthen}
\usepackage{tikz}
\usetikzlibrary{arrows.meta}
\usepackage{tkz-tab}
\usepackage{eurosym}

% OPTION PROF/ÉLÈVE
% Décommentez UNE des deux lignes suivantes :
\newboolean{prof}\setboolean{prof}{true}   % Version professeur (avec corrections des devoirs)
\newboolean{prof}\setboolean{prof}{true}  % Version élève (sans corrections des devoirs)

\definecolor{Couleur1}{HTML}{4A5D7A} %Th
\definecolor{Couleur2}{HTML}{A5B5D6} %Prop
\definecolor{Couleur3}{HTML}{A8C68A} %Méthode
\definecolor{Couleur6}{HTML}{8A9CC1} %Def
\definecolor{Couleur7}{HTML}{E6B459} %Rq Voc
\definecolor{Couleur8}{HTML}{D36740} %Titres
\definecolor{correction}{HTML}{D36740} % Couleur pour les corrections
\definecolor{coopmathscorr}{HTML}{F15929} % Couleur corrections coopmaths

% Configuration des marges
\geometry{
	top=2cm,
	bottom=2cm,
	inner=1.5cm,
	outer=1.5cm,
}

% Police
\renewcommand{\familydefault}{\sfdefault}

% Compteurs
\newcounter{seancenum}
\newcounter{seancenumD}
\setcounter{seancenum}{1}
\setcounter{seancenumD}{1}

% Configuration des en-têtes et pieds de page
\pagestyle{fancy}
\fancyhf{}
\ifthenelse{\boolean{prof}}{
    \fancyhead[L]{\textcolor{Couleur8}{\textbf{Corrections Automatismes - Classe de seconde - Version Professeur}}}
}{
    \fancyhead[L]{\textcolor{Couleur8}{\textbf{Corrections Devoirs - Séance 20 - Classe de seconde}}}
}
\fancyhead[R]{\textcolor{Couleur8}{\textbf{2025-2026}}}
\fancyfoot[L]{\textcolor{Couleur8}{Mme Bertail et Mme Pinto}}
\fancyfoot[R]{\textcolor{Couleur8}{\thepage}}
\renewcommand{\headrulewidth}{1pt}
\renewcommand{\headrule}{\hbox to\headwidth{\color{Couleur8}\leaders\hrule height \headrulewidth\hfill}}

% Environnements personnalisés
\newtcolorbox{seance}[1]{
    colback=Couleur6!10,
    colframe=Couleur1!65,
    fonttitle=\bfseries\large,
    title=Correction Séance \theseancenum #1,
    left=5mm,
    right=5mm,
    top=3mm,
    bottom=3mm,
    arc=0mm,
    after=\stepcounter{seancenum}
}

\newtcolorbox{seanceD}[1]{
    colback=Couleur6!5,
    colframe=Couleur6!45,
    fonttitle=\bfseries\large,
    title=Correction Devoirs - Séance \theseancenumD #1,
    left=5mm,
    right=5mm,
    top=3mm,
    bottom=3mm,
    arc=0mm,
    after=\stepcounter{seancenumD}
}

\begin{document}

\setcounter{seancenumD}{20}
\newpage
\begin{seanceD}{} %20

\begin{enumerate}
\item Développons l'expression :\\
 $\begin{aligned}A &=\left(-\dfrac{6}{7}y-4\right)^2-\left(\dfrac{4}{7}y+3\right)^2\\&=\left({\color{red}\boldsymbol{\left( \dfrac{-6y}{7}\right) ^2}}+{\color{blue}\boldsymbol{2\times \left( \dfrac{-6y}{7}\right) \times \left( -4\right) }}+{\color[HTML]{008002}\boldsymbol{\left( -4\right) ^2}}\right)-\left({\color{red}\boldsymbol{\left( \dfrac{4y}{7}\right) ^2}}+{\color{blue}\boldsymbol{2\times \left( \dfrac{4y}{7}\right) \times 3}}+{\color[HTML]{008002}\boldsymbol{3^2}}\right)\\&={\color{red}\boldsymbol{\frac{36y^2}{49}}}{\color{blue}\boldsymbol{+\frac{48y}{7}}}{\color[HTML]{008002}\boldsymbol{+16}}{\color{red}\boldsymbol{+\frac{-16y^2}{49}}}{\color{blue}\boldsymbol{+\frac{-24y}{7}}}{\color[HTML]{008002}\boldsymbol{-9}}\\&=({\color{red}\boldsymbol{\frac{36y^2}{49}+\frac{-16y^2}{49}}})+({\color{blue}\boldsymbol{\frac{48y}{7}+\frac{-24y}{7}}})+({\color[HTML]{008002}\boldsymbol{16-9}})\\&={\color[HTML]{f15929}\boldsymbol{\dfrac{20}{49}y^2+\dfrac{24}{7}y+7}}\end{aligned}$

\item  ${\color[HTML]{f15929}\boldsymbol{10^{-3}}} \leqslant 0{,}009\,4 \leqslant {\color[HTML]{f15929}\boldsymbol{10^{-2}}}$ car $10^{-3} = 0{,}001$ et $10^{-2} = 0{,}01.$

\item  $\overrightarrow{u}=\overrightarrow{BA}$ $\Leftrightarrow\begin{cases}-7=7-x_B\\-3=-5-y_B\end{cases}$ $\Leftrightarrow\begin{cases}x_B=7-(-7)\\y_B=-5-(-3)\end{cases}$ $\Leftrightarrow\begin{cases}x_B=14\\y_B=-2\end{cases}$, soit $B({\color[HTML]{f15929}\boldsymbol{14}}\,;\,{\color[HTML]{f15929}\boldsymbol{-2}})$.

\item  Comme $f(0)=10$, la fonction $f(x)=ax+b$ vérifie $a\times 0 + b = b = 10$ et par suite $f(x)=ax+10$.\\Comme $f(7)=-4$, le coefficient $a$ tel que de $f(x)=ax+10$ vérifie $a\times 7+10 = -4$ soit $7a=-14$.\\On en déduit $a=\dfrac{-14}{7}=-2$ et ainsi que $f(x)=-2x+10$.\\Posons $b$ l'antécédent de $-2$, alors $f(b)=-2\times b+10=-2$.\\On en déduit $-2b=-2-10=-12$.\\Donc $b=\dfrac{-12}{-2}=$ ${\color[HTML]{f15929}\boldsymbol{6}}$.

\item \begin{enumerate}[itemsep=1em]
	\item L'équation est de la forme $x^3=k$ avec $k=-27$. \\
              Le nombre dont le cube est $-27$ est $-3$ car $(-3)^3=-27$.\\
              Ainsi,   $S=\{-3\}$.
              
	\item L'équation est de la forme $x^2=k$ avec $k=121$. Comme  $121>0$ alors l'équation admet deux solutions : $-\sqrt{121}$ et $\sqrt{121}$.\\
                Comme $-\sqrt{121}=-11$ et $\sqrt{121}=11$ alors
                les solutions de l'équation peuvent s'écrire plus simplement : $-11$ et $11$.\\
                Ainsi,  $S={\color[HTML]{f15929}\boldsymbol{\{-11\,;\,11\}}}$.
\end{enumerate}

\end{enumerate}

\end{seanceD}

\end{document}